\section{Introduction}

Many alignment and oversight proposals begin from a dyadic picture: a human or institution supervises one powerful AI system. That framing remains useful for a range of problems, but it becomes incomplete once advanced AI deployment is expected to involve many interacting systems with partly aligned and partly conflicting objectives. Recent arguments about agentic AI emphasize this plural structure directly: intelligence may scale through interacting constellations of agents, services, and institutions rather than through a single monolithic system \citep{evans2025agentic}. On that view, alignment becomes partly an institutional problem. The question is not only whether one agent is well behaved, but whether the surrounding rules, protocols, and governance arrangements remain robust as strategic pressure increases.

This paper develops that framing through sequential commons. Commons settings provide a useful bridge between institutional theory and multi-agent learning because they make the central tension legible: individually beneficial behavior can still damage a shared system unless the rules of interaction are well designed \citep{gordon1954common,scott1955fishery,ostrom1990governing}. Work on sequential social dilemmas and multi-agent evaluation already points in this direction by showing that cooperation depends on temporally extended interaction, partner composition, and held-out social situations rather than only on one-shot payoffs \citep{leibo2017multiagent,meltingpot2021,socialjax2025}. Cooperative-AI work strengthens the same point from a game-theoretic angle: individually aligned systems do not guarantee collectively good outcomes, so the interaction structure matters in its own right \citep{conitzer2023foundations}.

The contribution here is therefore framed as \emph{institutional robustness under capability escalation}. In the benchmark, capability escalation is operationalized as repeated adversarial strategy injection: weaker strategies are replaced over generations, and stronger injectors search for more exploitative children. Governance conditions are then interpreted as candidate institutional designs. Under this translation, top-down controls, bottom-up coordination, and hybrid arrangements are not merely environment toggles. They are alternative ways of organizing restraint, intervention, and response inside a plural system.

The empirical program remains grounded in two linked studies, but their roles become clearer in this framing. Fishery Commons is the controlled precursor study. It asks which central governance signals survive repeated strategic turnover in a minimal renewable commons. Harvest Commons is the main study. It asks which governance architecture remains robust once local interaction, communication, and local side-payments are allowed to matter. The immediate next step is to make that second study more explicit by restoring the full architecture spectrum, organizing attacker strength into a capability ladder, and decomposing the welfare cost of robust governance by agent type and exposure.
